\documentclass[11pt]{article}

\usepackage[a4paper, margin=2.5cm]{geometry}
\usepackage{amsmath, amssymb}

\newcommand{\N}{\mathbb{N}}
\newcommand{\Pow}{\mathcal{P}}

\begin{document}

Decidan si la diferencia simétrica es asociativa o no:
\[
A\triangle(B\triangle C)=(A\triangle B)\triangle C.
\]
Si es asociativa, demuéstrelo. Si no, dé un contraejemplo.

¿Se les ocurre una manera fácil de describir $A\triangle(B\triangle C)$? ¿Funciona para $(A\triangle D)\triangle(B\triangle C)$?

\bigskip

La diferencia simétrica sí es asociativa:
\[
\boxed{A\triangle(B\triangle C)=(A\triangle B)\triangle C}.
\]
Por definición, $x\in X\triangle Y$ si y sólo si $x$ pertenece exactamente a
uno de $X$ y $Y$. Defino
\[
E:=\{x:x\text{ pertenece a un número impar de los conjuntos }A,B,C\}.
\]
Para tres conjuntos, esto significa que $x$ pertenece exactamente a uno de ellos o a los tres. Escrito de forma explícita,
\begin{align*}
E={}&[A\setminus(B\cup C)]\cup[B\setminus(A\cup C)]
\cup[C\setminus(A\cup B)]\\
&\cup(A\cap B\cap C).
\end{align*}

Voy a demostrar que ambos lados son iguales a $E$.

Primero, sea $x\in A\triangle(B\triangle C)$. Por definición de diferencia
simétrica ocurre exactamente uno de los siguientes casos:
\begin{enumerate}
  \item $x\in A$ y $x\notin B\triangle C$;
  \item $x\notin A$ y $x\in B\triangle C$.
\end{enumerate}
En el primer caso, $x$ pertenece a $A$ y pertenece a cero o a dos de los
conjuntos $B,C$; en total pertenece a uno o a tres de $A,B,C$. En el segundo
caso, no pertenece a $A$ y pertenece exactamente a uno de $B,C$; en total
pertenece a uno de los tres conjuntos. En cualquier caso, $x\in E$. Esto prueba
que
\[
A\triangle(B\triangle C)\subseteq E.
\]

Para la inclusión contraria, sea $x\in E$. Si $x\in A$, entonces el número de
pertenencias de $x$ entre $B,C$ debe ser par: pertenece a ambos o a ninguno.
Por tanto, $x\notin B\triangle C$, y así
$x\in A\triangle(B\triangle C)$. Si $x\notin A$, para que el número total de
pertenencias sea impar, $x$ debe pertenecer exactamente a uno de $B,C$; esto es,
$x\in B\triangle C$. De nuevo, $x\in A\triangle(B\triangle C)$. Así,
\[
E\subseteq A\triangle(B\triangle C),
\]
y por doble inclusión,
\[
A\triangle(B\triangle C)=E.
\]

Ahora considero el otro lado. Un elemento $x$ pertenece a
$(A\triangle B)\triangle C$ si y sólo si ocurre uno de estos dos casos:
\begin{enumerate}
  \item $x\in A\triangle B$ y $x\notin C$;
  \item $x\notin A\triangle B$ y $x\in C$.
\end{enumerate}
En el primer caso, $x$ pertenece exactamente a uno de $A,B$ y no pertenece a
$C$, de modo que tiene una pertenencia en total. En el segundo caso, pertenece
a cero o a dos de $A,B$ y también pertenece a $C$, de modo que tiene una o tres
pertenencias en total. Recíprocamente, si $x$ pertenece a un número impar de
$A,B,C$, al separar la pertenencia a $C$ se obtiene necesariamente uno de esos
dos casos. Por consiguiente,
\[
(A\triangle B)\triangle C=E.
\]
Como ambos lados son iguales a $E$, por transitividad de la igualdad concluyo
que
\[
A\triangle(B\triangle C)=(A\triangle B)\triangle C.
\]

La misma descripción funciona con cuatro conjuntos:
\[
\boxed{(A\triangle D)\triangle(B\triangle C)
=\{x:x\text{ pertenece a un número impar de }A,B,C,D\}}.
\]
Aquí, ``un número impar'' significa exactamente uno o exactamente tres.
En efecto, cada diferencia simétrica registra la paridad de las pertenencias:
dos apariciones se cancelan y una aparición permanece. Por eso agrupar los
cuatro conjuntos en parejas no cambia la paridad total.

\end{document}
